\lecture{13}{2. December 2025}{Svingning af rumligt kontinuerte systemer (del II)}

\begin{problemwithimage}{p13_1}{13.1}
  Vi betragter transversalsvingning af den viste bjælke, der er rumligt kontinuert og uniform.
  Bjælken, der har længden $L$, tværsnitsarealet $A$, arealinertimomentet $I$, elasticitetsmodulet $E$ og massedensiteten $\rho$, dissiperer ikke energi.
  \begin{enumerate}[label=\alph*)]
    \item Angiv de gældende rangbetingelser i forhold til transversalsvingning.
    \item Bestem egenfrekvenserne associeret med bjælkens transversalsvingning.
    \item Bestem de tilhørende egensvingningsformer.
  \end{enumerate}
\end{problemwithimage}

\begin{problem}{13.2}
  Bevis at bevægelsesligningen for vridningsegensvingning af en to-dimensionel, uniform bjælke er
  \begin{equation*}
    \frac{\partial ^2 \theta}{\partial t^2}  = \frac{G}{\rho} \frac{\partial ^2 \theta}{\partial x^2} 
  \end{equation*}
  hvori $\theta = \theta(x,t)$ er vridningen, $G$ er forskydningsmodulet, og $\rho$ er massedensiteten.
  Hint: Lad $T = T(x,t)$ være virdningsmomentet og $I_p$ det polære inertimoment, så gælder $T = GI_p \frac{\partial \theta}{\partial x} $.
\end{problem}

\begin{problemwithimage}{p13_3}{13.3}
  Vi betragter vridningssvingning af den viste bjælke, der er rumligt kontinuert og uniform.
  Bjælken, der har længden $L$, arealet $A$, det polære inertimoment $I_p$, forskydningsmodulet $G$ og massedensiteten $\rho$, dissiperer ikke energi.

  \begin{enumerate}[label=\alph*)]
    \item Angiv de gældende randbetingelser i forhold til vridningssvingning.
    \item Bestem egenfrekvenserne associeret med bjælkens vridningssvingning.
    \item Bestem de tilhørende egensvingningsformer.
  \end{enumerate}
\end{problemwithimage}
